\documentclass[a4paper]{article}
% Español (México) con XeLaTeX: fontspec para fuentes Unicode, polyglossia
% para la división silábica y los nombres del documento en la variante mexicana.
\usepackage{fontspec}
\usepackage{polyglossia}
\setdefaultlanguage[variant=mexican]{spanish}
% Los nombres chinos van como «pinyin (original)»: xeCJK da a los caracteres
% Han su propia fuente; sin ella XeLaTeX los omite con un aviso.
% FandolSong no cubre algunos caracteres de nombres (U+73FA); WenQuanYi Zen Hei sí.
\usepackage[AutoFallBack=true]{xeCJK}
\setCJKmainfont{FandolSong-Regular.otf}
\setCJKfallbackfamilyfont{\CJKrmdefault}{WenQuanYi Zen Hei}
% polyglossia no traduce los nombres de listings.
\AtBeginDocument{\providecommand{\lstlistingname}{}\renewcommand{\lstlistingname}{Listado}%
  \providecommand{\lstlistlistingname}{}\renewcommand{\lstlistlistingname}{Índice de listados}}
\usepackage{amsmath,amssymb}
\usepackage{graphicx}
\usepackage[margin=2.35cm]{geometry}
\usepackage[most]{tcolorbox}
\usepackage{etoolbox}
\usepackage{listings}
\usepackage{xcolor}
\usepackage{hyperref}
\usepackage{booktabs,longtable,tabularx,array}
\usepackage{subcaption}
\usepackage{float}
\usepackage{enumitem}
\usepackage{microtype}
\usepackage{fancyhdr}
\usepackage{needspace}

\definecolor{kbBack}{HTML}{F2F6FA}
\definecolor{kbFrame}{HTML}{9DB4CC}
\definecolor{kbTitle}{HTML}{52708F}
\definecolor{ibBack}{HTML}{FBF5E7}
\definecolor{ibFrame}{HTML}{C9A961}
\definecolor{ibTitle}{HTML}{7D5F1E}
\definecolor{wbBack}{HTML}{FCF2F0}
\definecolor{wbFrame}{HTML}{D6A096}
\definecolor{wbTitle}{HTML}{9E4F43}
\definecolor{dbBack}{HTML}{F4F7F3}
\definecolor{dbFrame}{HTML}{AEC2AC}
\definecolor{dbTitle}{HTML}{5F7F64}

\tcbset{softbox/.style={
  enhanced,breakable,boxrule=0.8pt,arc=2pt,left=3mm,right=3mm,
  top=2mm,bottom=2mm,coltitle=white,fonttitle=\bfseries,
  toptitle=0.6mm,bottomtitle=0.6mm
}}
\newtcolorbox{knowledgebox}[1]{softbox,colback=kbBack,colframe=kbFrame,colbacktitle=kbTitle,title={#1}}
\newtcolorbox{importantbox}[1]{softbox,colback=ibBack,colframe=ibFrame,colbacktitle=ibTitle,title={#1}}
\newtcolorbox{warningbox}[1]{softbox,colback=wbBack,colframe=wbFrame,colbacktitle=wbTitle,title={#1}}
\newtcolorbox{dialoguebox}[1]{softbox,colback=dbBack,colframe=dbFrame,colbacktitle=dbTitle,title={#1}}

\lstset{
  language=Python,basicstyle=\ttfamily\small,
  keywordstyle=\color{kbTitle}\bfseries,stringstyle=\color{wbTitle},
  commentstyle=\color{dbTitle},breaklines=true,frame=single,numbers=left,
  numberstyle=\tiny\color{gray},captionpos=b,columns=fullflexible,
  keepspaces=true,showstringspaces=false,extendedchars=true
}
\BeforeBeginEnvironment{lstlisting}{\Needspace{12\baselineskip}}

\hypersetup{colorlinks=true,linkcolor=kbTitle,urlcolor=kbTitle,citecolor=wbTitle}
\setlist{itemsep=0.25em,topsep=0.4em}
\setlength{\parindent}{2em}
\setlength{\parskip}{0.25em}
\newcolumntype{L}[1]{>{\raggedright\arraybackslash}p{#1}}

\providecommand{\notetitle}{Notas del curso: [tema del curso]}
\providecommand{\noteauthors}{Elaboradas a partir de los materiales públicos de Stanford CS25}
\providecommand{\notedate}{Versión reescrita en 2026}
\providecommand{\videochannel}{Stanford Online}
\providecommand{\videopublishdate}{2022}
\providecommand{\videoduration}{[duración del video]}
\providecommand{\videourl}{}
\providecommand{\videocoverpath}{cover.jpg}
\providecommand{\coursesite}{https://web.stanford.edu/class/cs25/}
\providecommand{\repourl}{https://github.com/hqhq1025/ai-course-notes}
\providecommand{\skillversion}{wdkns/wdkns-skills@39f1a04}

\newcommand{\figsource}[1]{\par\vspace{-0.3em}{\footnotesize\color{gray}\emph{Fuente: #1}}\par}
\newcommand{\teachervoice}[1]{\begin{knowledgebox}{Nota de clase}#1\end{knowledgebox}}
\newcommand{\term}[1]{\textbf{#1}}

\pagestyle{fancy}
\fancyhf{}
\fancyfoot[C]{\thepage}
\fancyfoot[R]{\footnotesize\color{gray}\href{\repourl}{AI Course Notes}}
\renewcommand{\headrulewidth}{0pt}
\renewcommand{\footrulewidth}{0pt}

\newcommand{\makecscover}{
\begin{titlepage}
\centering
\vspace*{0.7cm}
{\Large Stanford CS25 · Transformers United\par}
\vspace{0.8cm}
{\huge\bfseries \notetitle\par}
\vspace{0.6cm}
{\large \noteauthors\par}
\vspace{0.25cm}
{\large \notedate\par}
\vspace{0.9cm}
\ifdefempty{\videocoverpath}{}{
  \includegraphics[width=0.86\textwidth,height=0.43\textheight,keepaspectratio]{\videocoverpath}\par
}
\vfill
\begin{tcolorbox}[width=0.92\textwidth,colback=black!2!white,colframe=black!45,boxrule=0.7pt,arc=2pt]
\textbf{Autor o canal del video}: \videochannel\par
\textbf{Fecha de publicación}: \videopublishdate\par
\textbf{Duración del video}: \videoduration\par
\textbf{Sitio del curso}: \href{\coursesite}{\nolinkurl{\coursesite}}\par
\ifdefempty{\videourl}{}{\textbf{Enlace del video}: \href{\videourl}{\nolinkurl{\videourl}}\par}
\textbf{Normas de elaboración}: \skillversion\ + las normas source-first / teacher-voice / visual-QA de este repositorio
\end{tcolorbox}
\end{titlepage}
\tableofcontents
\newpage
}


\renewcommand{\notetitle}{Lecture 18: Neuroscience-Inspired Artificial Intelligence}
\renewcommand{\noteauthors}{Elaborado a partir de la clase de Trenton Bricken y Will Dorrell, los subtítulos oficiales hechos por personas, las diapositivas recuperadas y los artículos indicados}
\renewcommand{\notedate}{CS25 V2 · versión reescrita source-first de 2026}
\renewcommand{\videochannel}{Stanford Online}
\renewcommand{\videopublishdate}{Clase: 2023-03-07; publicación: 2023-09-01}
\renewcommand{\videoduration}{1:22:13}
\renewcommand{\videourl}{https://www.youtube.com/watch?v=L4DC7e6g2iI}
\renewcommand{\videocoverpath}{cover.jpg}

\newcommand{\lecturefigure}[3]{%
\begin{figure}[H]
  \centering
  \includegraphics[width=0.94\textwidth,height=0.54\textheight,keepaspectratio]{slides-images/#1}
  \caption{#2}
  \figsource{#3}
\end{figure}
}

\let\standardtableofcontents\tableofcontents
\renewcommand{\tableofcontents}{%
  \begingroup
  \small
  \setlength{\parskip}{0pt}
  \standardtableofcontents
  \endgroup
}

\begin{document}
\makecscover

\section{Auditoría de fuentes: una clase, dos rutas desde la neurociencia}

Esta clase no es una conferencia sobre un solo artículo: dos ponentes parten de regiones cerebrales distintas y de niveles de cómputo distintos, y cada uno llega por su lado a la Transformer-like computation. En los primeros 42 minutos, Trenton Bricken parte de la geometría de alta dimensión de la Sparse Distributed Memory (SDM, memoria distribuida dispersa), explica por qué la hypersphere intersection produce pesos exponenciales que aproximan una softmax y propone que el cerebellum (cerebelo) podría implementar un circuito de lectura y escritura similar. En los 39 minutos siguientes, Will Dorrell parte del cognitive map, el hippocampus y la entorhinal cortex, presenta cómo la Tolman--Eichenbaum Machine (TEM) separa la estructura del contenido y cómo la modern Hopfield memory la reformula como attention.

\subsection{Corrección del video y del alcance erróneos de la versión anterior}

El video \texttt{dEFn6nnoC-8} enlazado en la versión anterior es la defensa doctoral de Trenton Bricken, realizada el 31 de marzo de 2025 y publicada en mayo de 2025; no corresponde a esta clase. La grabación oficial correcta es \texttt{L4DC7e6g2iI} de Stanford Online, y la fecha de la clase es el 7 de marzo de 2023. Además, la versión anterior solo cubría al primer ponente y agregaba contenido que no aparece en la clase, como el enrutamiento de MoE, un dashboard de producción, un playbook de gobernanza, incident response y el despliegue intermodal; todos esos párrafos se eliminaron en esta reescritura.

\begin{longtable}{@{}L{0.18\textwidth}L{0.29\textwidth}L{0.41\textwidth}@{}}
\toprule
Nivel de evidencia & Material local & Función didáctica \\
\midrule
Grabación oficial & 1:22:13, dos ponentes, 66 estados didácticos finales & Determina el orden de la clase, las fórmulas, las derivaciones manuscritas, las preguntas y respuestas y las conclusiones \\
Subtítulos oficiales hechos por personas & 1,950 caption y timed transcript & Recuperan la motivación, los límites de las analogías, el alcance de los experimentos y los detalles de neurociencia \\
Artículo de la primera parte & Attention Approximates SDM & Verificar la intersección en alta dimensión, la aproximación exponencial, la SDM continua y el learned coefficient \\
Artículo de la segunda parte & Artículos sobre TEM y el hippocampal Transformer & Verificar la descomposición estructura/contenido, la actualización de la memoria, el grid-like code y su correspondencia con attention \\
Auditoría visual & 40 slide estándar + 26 estados finales manuscritos & Cubrir por completo ambas partes de la clase sin convertir cada trazo escrito en una captura repetida \\
\bottomrule
\end{longtable}

\begin{warningbox}{Límites históricos y de la evidencia}
Esta clase reconstruye la sesión del 2023-03-07. Una correspondencia matemática indica que «cierto cómputo puede implementarse de forma aproximada mediante otra forma»; no demuestra que el cerebro ejecute un Transformer entrenado. Una asignación a regiones cerebrales indica que «el circuito tiene componentes candidatos para implementarlo»; no significa que se haya confirmado experimentalmente. La semejanza entre TEM y el Transformer se da en el nivel de la recuperación de memoria, el estado posicional y la generalización estructural; tampoco es una identidad anatómica neurona por neurona.
\end{warningbox}

\subsection{Pregunta de investigación y ruta de la primera parte}

La portada presenta el trabajo de Bricken y Cengiz Pehlevan; la diapositiva de motivación resume el problema en una frase: la attention del Transformer es muy eficaz, pero la softmax es una heuristic; ¿podría surgir de forma natural a partir de propiedades simples de los vectores de alta dimensión? La diapositiva de Agenda mantiene el orden: primero la intuición visual, después la derivación matemática y luego la discusión de los componentes del Transformer y de su plausibilidad biológica. Este orden importa, porque si la intersección de esferas en alta dimensión se escribe directamente como una función especial, se oculta su intuición como memoria: el «número de direcciones que se traslapan».

\lecturefigure{slide-01-title-attention-sdm.jpg}{Primera parte: Attention Approximates Sparse Distributed Memory}{Diapositiva recuperada de la grabación oficial V001; 00:00:22}

\lecturefigure{slide-02-why-care.jpg}{Por qué vale la pena estudiarlo: de la heuristic softmax a la geometría de la memoria en alta dimensión}{Diapositiva recuperada de la grabación oficial V002; 00:01:20}

\lecturefigure{slide-03-presentation-overview.jpg}{Ruta de la primera parte: SDM, attention, su correspondencia y la biological plausibility}{Diapositiva recuperada de la grabación oficial V003; 00:01:52}

\teachervoice{Bricken dice explícitamente que primero dará una high-level visual intuition y después entrará en las matemáticas, y anima a los estudiantes a interrumpir con preguntas en cualquier momento. Por eso estas notas siguen el mismo orden: primero se dibujan con claridad la escritura, la lectura y la intersección, y después se derivan los pesos exponenciales. No se trata de evitar las fórmulas, sino de que cada símbolo tenga primero un significado operativo.}

\subsection{Las tres preguntas que comparten ambas partes}

Aunque SDM y TEM provienen de tradiciones teóricas distintas, las dos partes indagan tres niveles de relación: en el \textbf{nivel algorítmico}, cómo recupera el modelo los recuerdos pertinentes; en el \textbf{nivel de representación}, cómo se codifican la estructura y el contenido; en el \textbf{nivel de implementación}, qué circuitos neuronales o matrices aprendibles pueden ejecutar el cómputo. El valor de la clase no está en reducir los tres niveles a la frase «el cerebro es un Transformer», sino en comparar cuándo una correspondencia entre niveles es exacta y cuándo es solo una inspiración.

\begin{importantbox}{Escala de evidencia para leer toda la clase}
Distinga, en orden: equivalencia o aproximación de fórmulas, resultados empíricos en modelos entrenados, similitud de representaciones neuronales, posibilidad de implementación en circuitos candidatos y evidencia causal en el cerebro real. Cuanto más adelante en la escala, más exigente es la evidencia requerida. Una fórmula de attention elegante no demuestra por sí sola la función de una región cerebral, y un grid pattern similar tampoco demuestra por sí solo que dos sistemas usen el mismo algoritmo.
\end{importantbox}

\subsection{Resumen de la sección}

La fuente correcta de la clase incluye a dos ponentes y dos rutas complementarias. La primera explica los pesos de attention desde la memoria asociativa de alta dimensión; la segunda explica la generalización estructural y la attention-like retrieval desde los mapas cognitivos y la hippocampal memory. Todo el contenido posterior debe mantener los niveles de evidencia entre el algoritmo, la representación y la implementación biológica.
\section{Sparse Distributed Memory: cómo escribir y cómo leer a partir de consultas con ruido}

La primera parte de la clase todavía no trata el Transformer; construye un sistema de memoria asociativa. El objetivo no es leer un registro con una dirección exacta, sino que una query cercana siga recuperando el pattern correcto, y obtener al mismo tiempo alta capacidad, tolerancia a fallos y almacenamiento distribuido. La técnica central de SDM es fijar un conjunto de hard locations dispersas, de modo que cada escritura y cada lectura actúen sobre un grupo de posiciones dentro del radio de la query.

\subsection{Objetivos de diseño: capacidad, robustez y fault tolerance}

La diapositiva de requisitos enumera cuatro: high memory capacity, robustness to query noise, biological plausibility y fault tolerance. La siguiente diapositiva separa sparse de distributed: el espacio de direcciones tiene una dimensión altísima y las neuronas reales ocupan solo una fracción mínima de las posiciones posibles; sin embargo, un pattern activa a la vez varias posiciones cercanas, por lo que ninguna dirección individual contiene la memoria completa.

\lecturefigure{slide-04-brain-memory-requirements.jpg}{Requisitos de diseño de una memoria asociativa de tipo cerebral}{Diapositiva recuperada del video oficial V004; 00:02:12}

\lecturefigure{slide-05-sdm-unique.jpg}{Los dos sentidos de SDM: sparse y distributed}{Diapositiva recuperada del video oficial V005; 00:02:48}

\begin{knowledgebox}{Glosario: address, pattern, value y hard location}
\term{pattern} es el vector de alta dimensión que se escribe o que se usa como query; \term{hard location} es una dirección dispersa predefinida por el sistema; \term{address/key} determina qué posiciones se activan; \term{value} es el contenido que se escribe y que se devuelve en la consulta. En la auto-associative memory clásica, key y value pueden ser iguales; en la hetero-associative memory pueden ser distintos.
\end{knowledgebox}

Sea el espacio binario de direcciones $\{0,1\}^{n}$, la $i$-ésima hard location $h_i$, la dirección de escritura $p$ y el radio $d$. El indicador de activación es
$$
a_i(p)=\mathbf{1}\!\left[D_H(h_i,p)\le d\right],
$$
donde $D_H$ es la Hamming distance, $n$ es la dimensión de los vectores y $d$ controla cuántas posiciones se activan. En alta dimensión, la mayoría de los vectores aleatorios están muy lejos entre sí, y aun así un radio local puede contener suficientes hard locations para realizar un cómputo distribuido.

\subsection{Escritura: posiciones cercanas y superposition}

En la primera figura de escritura, el pattern verde activa las hard locations dentro de su radio; en la segunda se agrega el pattern naranja, y los radios de ambas escrituras pueden traslaparse; en la tercera se agrega el pattern azul, y las posiciones del traslape guardan a la vez las contribuciones de varios value. La superposition no concatena varias memorias: las acumula sobre el mismo vector de almacenamiento, y en la lectura el ruido se elimina mediante la agregación estadística de muchas posiciones.

\lecturefigure{slide-06-sdm-write-pattern1.jpg}{SDM escribe el primer pattern: se activan las hard locations dentro del radio}{Diapositiva recuperada del video oficial V006; 00:03:26}

\lecturefigure{slide-07-sdm-write-pattern2.jpg}{Escritura del segundo pattern: los vecindarios de las direcciones pueden traslaparse}{Diapositiva recuperada del video oficial V007; 00:03:42}

\lecturefigure{slide-08-sdm-write-superposition.jpg}{Varios pattern se almacenan superpuestos en las hard locations}{Diapositiva recuperada del video oficial V008; 00:04:14}

Si la $i$-ésima posición guarda un vector de contadores $w_i$, escribir el value $v$ en la dirección $p$ puede expresarse como
$$
w_i\leftarrow w_i+a_i(p)\,v.
$$
Aquí $v$ puede ser un vector bipolar o un vector de valores reales, y $a_i(p)$ determina si la posición participa. Una posición acumula varios $v$, y la capacidad y el ruido provienen del mismo mecanismo: cuantos más pattern comparten posiciones, mayor es el aprovechamiento del almacenamiento y también mayor la interferencia cruzada.

\begin{lstlisting}[caption={Escritura distribuida de SDM (pseudocódigo conceptual)}]
def sdm_write(hard_locations, counters, address, value, radius):
    for index, location in enumerate(hard_locations):
        if hamming(location, address) <= radius:
            counters[index] += value
\end{lstlisting}

\teachervoice{El ponente recuerda varias veces al público que primero piense en estas «neuron» como hard locations abstractas y que solo después se relacionarán con la biología. Este orden evita interpretar un diagrama bidimensional de círculos como un espacio real dentro del cerebro: la distancia en la figura representa una Hamming distance de alta dimensión, y los círculos son solo una proyección didáctica.}

\subsection{Lectura: de las neuronas activadas a la intersección de pattern}

La operación de lectura tiene dos perspectivas equivalentes. En la neuron view, la query activa las posiciones cercanas, se suman los vectores almacenados en cada posición y luego se aplica majority vote dimensión por dimensión; la pattern view abstrae las hard locations y examina directamente el tamaño de la intersección entre la query sphere y cada write sphere. Cuanto mayor es la intersección, más posiciones se activaron tanto al escribir ese pattern como al leer con la query actual, y mayor es el peso del value de ese pattern.

\lecturefigure{slide-09-sdm-read-neurons.jpg}{Neuron view: una query con ruido activa posiciones cercanas y agrega lo almacenado}{Diapositiva recuperada del video oficial V009; 00:05:02}

\lecturefigure{slide-10-sdm-read-patterns.jpg}{Pattern view: la sphere intersection determina el peso de cada memoria}{Diapositiva recuperada del video oficial V010; 00:05:40}

\lecturefigure{slide-11-circle-intersection.jpg}{Intersección y distancia entre dos esferas de alta dimensión}{Diapositiva recuperada del video oficial V011; 00:05:56}

La forma continua de la lectura puede escribirse como
$$
r(q)=\sum_i a_i(q)w_i,
\qquad
\hat v_j=g(r_j)=\mathbf{1}[r_j>0],
$$
donde $q$ es la query, $r(q)$ es la lectura total de las posiciones activadas y $g$ es un umbral por dimensión o un majority decoder. Al sustituir la fórmula de escritura, la contribución de cada pattern $p_\mu$ a la salida es proporcional a
$$
I\!\left(D_H(p_\mu,q);d,n\right)
=\left|B(p_\mu,d)\cap B(q,d)\right|,
$$
es decir, al tamaño de la intersección de dos Hamming ball; $B(p,d)$ denota la esfera con centro en $p$ y radio $d$.

\lecturefigure{slide-12-query-read-equation.jpg}{Las cuatro etapas de la fórmula de query read: intersección, ponderación, normalización y decodificación}{Diapositiva recuperada del video oficial V012; 00:07:18}

\begin{knowledgebox}{Lectura de la figura: por qué una noisy query aún puede recuperar el pattern}
La query no tiene que ser igual a la dirección original: mientras siga traslapándose ampliamente con la sphere del pattern correcto, activará muchas de las posiciones donde se escribió ese value. Cuando errores aleatorios alteran unos cuantos bit, la intersección correcta suele seguir siendo mayor que la de los pattern no relacionados; el majority vote promedia el ruido independiente. La recuperación falla cuando los pattern son demasiado densos, el radio es demasiado grande o la query está demasiado lejos del pattern original.
\end{knowledgebox}

\subsection{Resumen de la sección}

SDM usa direcciones dispersas de alta dimensión para realizar escrituras distribuidas y la sphere intersection para el direccionamiento por contenido. El radio de escritura controla a la vez la capacidad y la interferencia; la lectura elimina el ruido mediante la agregación de las posiciones activadas y la majority decoding. En la siguiente sección basta demostrar que el intersection weight tiene una forma exponencial similar a softmax para conectar estas operaciones de memoria con attention.
\section{Transformer Attention: primero, definir con claridad qué se compara}

La clase no aplica la SDM directamente al Transformer; antes repasa la attention de principio a fin. El propósito es fijar qué se va a comparar: cómo la similitud entre query y key produce los pesos, qué hace softmax, por qué value es distinto de key y cómo se escribe en forma matricial el vector de actualización final.

\subsection{Semántica operativa de Q, K y V}

Esta subsección fija primero la semántica operativa de Q, K y V, de la que dependen las comparaciones geométricas posteriores. Después de la diapositiva de transición, el ejemplo de lenguaje trata la posición que se va a predecir como query, usa la key de cada token del contexto para el emparejamiento y trata la value como el contenido que realmente se transfiere a la salida. La diapositiva de Q/K/V divide el proceso en generación de proyecciones, comparación query--key, normalización y ponderación de las values; la diapositiva siguiente muestra cómo los pesos normalizados multiplican cada value y luego se suman.

\lecturefigure{slide-13-transformer-attention-overview.jpg}{Transición a la sección de Transformer attention}{Diapositiva V013 recuperada de la grabación oficial; 00:07:20}

\lecturefigure{slide-14-attention-example.jpg}{Ejemplo de attention en un modelo de lenguaje}{Diapositiva V014 recuperada de la grabación oficial; 00:08:02}

\lecturefigure{slide-15-qkv-steps.jpg}{Los cuatro pasos de cálculo con keys, queries y values}{Diapositiva V015 recuperada de la grabación oficial; 00:08:40}

\lecturefigure{slide-16-weighted-value-sum.jpg}{Cómo los pesos de softmax agregan los value vectors}{Diapositiva V016 recuperada de la grabación oficial; 00:09:00}

\begin{importantbox}{Key y value no deben confundirse}
La key responde «qué tan relacionado está este recuerdo con la query», y la value responde «si está relacionado, qué debe recuperarse». Al separarlas, el modelo puede emparejar por sintaxis o por identidad de entidad y, aun así, devolver características distintas que sean útiles para el siguiente token; esto también corresponde a la hetero-association de la SDM, en la que address y stored content pueden separarse.
\end{importantbox}

\subsection{Ejemplo ficticio y ecuación de actualización completa}

La diapositiva Fictional example usa tokens como ``cat'' y ``sat'' para mostrar que una key muy relacionada puede devolver una value que contiene información semántica y sintáctica. A continuación, la diapositiva Full query update escribe juntas la proyección y la actualización. Sea $X$ la matriz de tokens de entrada, $x_q$ el vector de la posición de consulta y $W_Q,W_K,W_V$ las matrices de proyección; entonces
$$
q=W_Qx_q,\quad K=XW_K,\quad V=XW_V,
$$
$$
\operatorname{Attn}(q,K,V)
=V^{\top}\operatorname{softmax}\!\left(\beta Kq\right).
$$
Aquí $q$ es la query, cada fila de $K$ es una key, cada fila de $V$ es una value y $\beta$ es el coeficiente de temperatura o escala. Si se usa scaled dot-product, lo habitual es $\beta=1/\sqrt{d_k}$.

\lecturefigure{slide-17-fictional-example.jpg}{Ejemplo ficticio: al emparejar una key se devuelven características de value distintas}{Diapositiva V017 recuperada de la grabación oficial; 00:10:14}

\lecturefigure{slide-18-full-query-update.jpg}{Expresión matricial de la query update completa}{Diapositiva V018 recuperada de la grabación oficial; 00:10:30}

\subsection{Dot product, softmax y value summation}

Tres diapositivas amplían, cada una, una operación central de la attention. La diapositiva de dot product explica que la similitud proviene de $k_i^{\top}q$; la de softmax explica que la exponenciación amplía las diferencias entre puntuaciones altas y normaliza; la de summation explica que la salida es una combinación convexa de las values. Lo esencial de softmax no es la etiqueta de «producir probabilidades», sino que los pesos varían exponencialmente con la similitud, que es justamente la forma que la sección siguiente busca recuperar a partir de la intersección de esferas en alta dimensión.

\lecturefigure{slide-19-query-key-dot-product.jpg}{El dot product query--key forma los similarity logits}{Diapositiva V019 recuperada de la grabación oficial; 00:10:50}

\lecturefigure{slide-20-softmax.jpg}{Softmax amplifica los términos de alta similitud y completa la normalización}{Diapositiva V020 recuperada de la grabación oficial; 00:11:44}

\lecturefigure{slide-21-summation-over-values.jpg}{Suma de los value vectors según los pesos de atención}{Diapositiva V021 recuperada de la grabación oficial; 00:12:02}

\lecturefigure{slide-22-full-attention-equation.jpg}{Flujo completo de la attention y ecuación de actualización final}{Diapositiva V022 recuperada de la grabación oficial; 00:12:18}

Para la $i$-ésima key, el peso de softmax es
$$
\alpha_i(q)=\frac{\exp(\beta k_i^{\top}q)}
{\sum_j\exp(\beta k_j^{\top}q)},
\qquad
y(q)=\sum_i\alpha_i(q)v_i.
$$
Aquí $\alpha_i$ es el normalized retrieval weight y $v_i$ es el contenido del $i$-ésimo recuerdo. Al compararla con la SDM, la pregunta se convierte en: si el sphere-intersection count puede escribirse como $\exp(\text{similarity})$ y si, tras normalizar, se obtiene el mismo $\alpha_i$.

\begin{lstlisting}[caption={Scaled dot-product attention (ilustración con una sola query)}]
def attention(query, keys, values, beta):
    logits = beta * (keys @ query)
    weights = softmax(logits)
    return weights @ values
\end{lstlisting}

\teachervoice{El ponente advierte en particular que «what you pay attention to» y «what you get out» pueden ser distintos: la query y las keys determinan la relevancia, y las values determinan el contenido recuperado. Esta idea es esencial para separar más adelante la cerebellar activation pathway de la write pathway, y también se conecta con la content/structure factorization de la segunda clase.}

\subsection{Resumen de la sección}

La attention es una memoria direccionable por contenido: la query y las keys producen similitudes, softmax las convierte en pesos normalizados exponencialmente y las values se recuperan ponderadas. Siempre que la sphere intersection de la SDM decaiga exponencialmente con la distancia y la distancia pueda reescribirse como producto punto, ambas operaciones de lectura comparten la misma estructura de cálculo.
\section{Por qué la intersección de esferas de alta dimensión aproxima Softmax}

Ahora entramos en la derivación central de la primera clase. La intuición es la siguiente: cuanto más alejados están los centros de dos esferas de alta dimensión con el mismo radio, menor es el volumen de su intersección; en el rango de parámetros relevante, esta disminución se aproxima a una función exponencial. Si los vectores están normalizados, la Hamming distance puede reescribirse como cosine/dot-product similarity, de modo que el intersection weight adopta la forma exponencial de un logit de softmax.

\subsection{De la curva de intersección al decaimiento exponencial}

La figura que sigue a la portada de la sección pone en el eje horizontal la Hamming distance entre los centros de las dos esferas y en el eje vertical el logaritmo del expected intersection size. Que los puntos se aproximen a una recta significa que el tamaño original de la intersección decae de forma aproximadamente exponencial:
$$
I(\delta;d,n)\approx c_1\exp(-c_2\delta),
$$
donde $\delta=D_H(p,q)$, $d$ es el radio de escritura/lectura, $n$ es la dimensión, y $c_1,c_2$ dependen de $n,d$. Esta expresión no es exacta para todas las distancias, sino que ofrece una aproximación válida en la región de similitud que la atención usa en la práctica.

\lecturefigure{slide-23-attention-approximates-sdm.jpg}{Pregunta central: cómo aproxima la attention del Transformer a SDM}{Diapositiva V023 recuperada de la grabación oficial; 00:12:26}

\lecturefigure{slide-24-exponential-circle-decay.jpg}{La sphere intersection de alta dimensión decae de forma aproximadamente exponencial con la distancia entre centros}{Diapositiva V024 recuperada de la grabación oficial; 00:12:46}

\begin{warningbox}{La figura de círculos en dos dimensiones solo aporta intuición}
En dos dimensiones, el área de intersección de dos círculos no muestra por sí sola el fenómeno de concentración propio de las altas dimensiones. La conclusión real depende de que $n$ sea grande y de que el radio y la dimensión guarden una proporción adecuada. No puede deducirse un teorema de que la figura «parezca exponencial»; el curso respalda la aproximación con gráficas numéricas y análisis asintótico a la vez.
\end{warningbox}

\subsection{Cómo la Hamming distance se convierte en cosine similarity}

Para vectores bipolar $a,b\in\{-1,+1\}^{n}$, cada bit igual aporta $+1$ al producto punto y cada bit distinto aporta $-1$; por lo tanto,
$$
D_H(a,b)=\frac{n-a^{\top}b}{2}.
$$
Si además se define $\hat a=a/\sqrt n$, $\hat b=b/\sqrt n$, entonces
$$
D_H(a,b)=\frac{n}{2}\left(1-\hat a^{\top}\hat b\right).
$$
Al sustituir esto en el decaimiento exponencial, el término constante se cancela con la normalización y el peso restante es proporcional a $\exp(\beta\hat a^{\top}\hat b)$. Este es el puente matemático entre la memoria asociativa basada en distancias y la dot-product attention.

\lecturefigure{slide-25-hamming-cosine-derivation.jpg}{De la Hamming distance al dot product normalizado}{Diapositiva V025 recuperada de la grabación oficial; 00:16:24}

\lecturefigure{slide-26-exponential-fit-regression.jpg}{Ajuste lineal sobre el log intersection size}{Diapositiva V026 recuperada de la grabación oficial; 00:17:18}

\begin{knowledgebox}{Lectura de la figura: qué respalda la gráfica de ajuste}
Al aumentar el eje horizontal, el log intersection size disminuye de forma aproximadamente lineal, lo que respalda la aproximación exponencial; la pendiente corresponde al effective beta y la ordenada al origen se cancela en la normalización de softmax. La figura no demuestra que todas las dimensiones y radios funcionen igual de bien; la cola curvada es justamente la región donde hay que revisar el error de aproximación.
\end{knowledgebox}

\subsection{Alinear término por término las dos notaciones}

La diapositiva Attention-in-SDM-notation alinea primero la query, la pattern address y el stored value; después, la diapositiva de la equivalence central sustituye el circle-intersection weight por un exponential dot product. Tras la normalización, la lectura en la SDM pattern view es
$$
\hat y(q)
=\frac{\sum_\mu I(D_H(p_\mu,q);d,n)v_\mu}
{\sum_\nu I(D_H(p_\nu,q);d,n)}
\approx
\sum_\mu
\frac{\exp(\beta p_\mu^{\top}q)}
{\sum_\nu\exp(\beta p_\nu^{\top}q)}v_\mu.
$$
Aquí $p_\mu$ es la $\mu$-ésima key/pattern, $v_\mu$ es el value, y $\beta$ depende conjuntamente de la dimensión, el radio, la norma de los vectores y el coeficiente de aproximación. El lado derecho es exactamente el attention read.

\lecturefigure{slide-27-attention-sdm-notation.jpg}{Reescritura de la attention con notación de SDM}{Diapositiva V027 recuperada de la grabación oficial; 00:17:20}

\lecturefigure{slide-28-sdm-attention-equivalence.jpg}{Correspondencia entre el SDM intersection weighting y la attention softmax}{Diapositiva V028 recuperada de la grabación oficial; 00:17:22}

\subsection{Experimentos numéricos y condiciones de validez}

La diapositiva Parameter experiments compara la intersection real con el ajuste exponencial para distintos valores de $n,d$; la diapositiva Mapping conditions subraya dos condiciones prácticas: los vectores deben estar aproximadamente normalizados en $L^2$, algo que suele garantizarse mediante LayerNorm o el norm behavior; y el effective coefficient debe coincidir con la intersection slope. Si las key/query norms varían libremente, también modifican la temperature, y no puede tomarse $1/\sqrt{d_k}$ como el único beta.

\lecturefigure{slide-29-parameter-experiments.jpg}{Experimentos de aproximación de la intersection con distintas dimensiones y radios}{Diapositiva V029 recuperada de la grabación oficial; 00:18:38}

\lecturefigure{slide-30-mapping-conditions.jpg}{Condiciones de normalization y de coefficient necesarias para que la correspondencia se cumpla}{Diapositiva V030 recuperada de la grabación oficial; 00:20:40}

\teachervoice{El ponente no oculta el «aproximadamente igual». Distingue una y otra vez entre close approximation y exact identity, y señala que la attention dispersa/dura, los vectores sin normalizar o un coefficient distinto se apartan de la correspondencia estándar. El método realmente transferible consiste en escribir primero las assumptions y luego preguntar si el modelo entrenado las cumple, en lugar de declarar la equivalencia en cuanto aparece una softmax.}

\begin{importantbox}{Versión mínima de la conclusión central}
En alta dimensión, con un radio adecuado, vectores normalizados y un coefficient que coincida, el número de elementos en la intersección de las esferas de lectura y escritura de SDM decae de forma aproximadamente exponencial con la distancia pattern--query; la distancia puede reescribirse como dot product, y tras la normalización se obtiene la softmax attention. Esta conclusión explica un origen geométrico de softmax, pero no es el único origen de softmax.
\end{importantbox}

\subsection{Resumen de la sección}

La intersection concentration en alta dimensión produce pesos exponenciales en función de la distancia, y la bipolar Hamming distance equivale a una transformación afín del dot product normalizado; por ello, el SDM pattern read aproxima la attention. Los experimentos numéricos respaldan esta aproximación, pero LayerNorm, la norma de los vectores, beta y la región de parámetros determinan la calidad de la correspondencia.
\section{SDM continua y explicación de los componentes del Transformer}

La SDM clásica usa binary address y hard radius; el Transformer usa continuous vector y softmax. A continuación, la clase vuelve continua la SDM y plantea una pregunta más abierta: si un Transformer entrenado cae realmente en un «intervalo de parámetros de SDM útil», y si la FFN, la LayerNorm y la memory capacity pueden integrarse en este mismo marco.

\subsection{Del binary radius al kernel continuo}

La diapositiva de Continuous SDM explica que no es necesario cuantizar la entrada en bits. Para vectores reales normalizados, basta con usar directamente un dot-product kernel exponencial para definir una lectura continua:
$$
\operatorname{CSDM}(q)=
\sum_\mu\operatorname{softmax}_\mu(\beta p_\mu^{\top}q)v_\mu.
$$
Esto tiene la misma forma que la attention y también puede implementarse con un perceptrón multicapa o con una associative-memory network. La versión continua conserva el principio de que «las direcciones cercanas contribuyen más», pero convierte el radio duro en pesos suaves.

\lecturefigure{slide-31-continuous-sdm.jpg}{Continuous SDM: del binary Hamming space a vectores reales}{Diapositiva V031 recuperada de la grabación oficial; 00:21:32}

\begin{knowledgebox}{Qué es la effective beta}
$\beta$ determina la retrieval sharpness: una $\beta$ pequeña hace que muchas memories participen promediadas; una $\beta$ grande hace que domine la memory más similar. En el Transformer, la scale explícita, las key/query norms y la LayerNorm modifican conjuntamente la temperatura efectiva, por lo que para medir el learned coefficient hay que observar los logits reales y no solo la constante de la fórmula.
\end{knowledgebox}

\subsection{¿Se parece la Attention entrenada a una SDM «buena»?}

La diapositiva de la pregunta propone una prueba inversa: si la attention implementa una SDM, ¿corresponden los parámetros entrenados a un radio/temperature con buena capacidad y buena robustez al ruido? La diapositiva de Learned beta muestra la distribución del coefficient en distintas heads, lo que indica que el modelo no tiene una temperatura uniforme. Cada head puede encargarse de tareas distintas, como copiar, sintaxis local o recuperación de largo alcance, por lo que la retrieval sharpness óptima también puede diferir.

\lecturefigure{slide-32-transformer-sdm-question.jpg}{Pregunta abierta: ¿corresponde un Transformer entrenado a parámetros de SDM óptimos?}{Diapositiva V032 recuperada de la grabación oficial; 00:22:18}

\lecturefigure{slide-33-learned-beta.jpg}{La effective beta aprendida por distintas attention heads}{Diapositiva V033 recuperada de la grabación oficial; 00:23:44}

\begin{warningbox}{No tomar la entropy como indicador monótono de calidad}
Una entropy alta puede indicar que hace falta integrar varias memorias, pero también que la recuperación no está enfocada; una entropy baja puede ser una recuperación precisa, pero también un exceso de confianza en un error. La perspectiva de la SDM ofrece una explicación en términos de capacidad/ruido, pero el juicio debe combinarse con la función de la head y la pérdida de la tarea.
\end{warningbox}

\subsection{FFN, LayerNorm y el memory framework}

La diapositiva de Transformer components integra la broader architecture en el marco de la SDM: la FFN puede verse como un conjunto de key--value memories de largo plazo; la LayerNorm ayuda a estabilizar la norma de los vectores y el coefficient; la attention es una pattern-based retrieval explícita. Por último, la diapositiva de summary presenta las conclusiones junto a las preguntas futuras y subraya que estas correspondencias pueden generar hipótesis de investigación, pero no constituyen una teoría unificada ya demostrada.

\lecturefigure{slide-34-transformer-components-converge.jpg}{Cómo convergen los componentes del Transformer en el SDM framework}{Diapositiva V034 recuperada de la grabación oficial; 00:25:06}

\lecturefigure{slide-35-sdm-summary.jpg}{Resumen de la parte teórica de la primera exposición y preguntas abiertas}{Diapositiva V035 recuperada de la grabación oficial; 00:26:20}

\begin{longtable}{@{}L{0.20\textwidth}L{0.31\textwidth}L{0.37\textwidth}@{}}
\toprule
Componente del Transformer & Perspectiva de SDM/memoria & Limitación que debe conservarse \\
\midrule
Attention & Pattern-based retrieval precisa en el contexto actual & Restringida por la context window; almacena todas las keys/values \\
FFN & Key--value associations de largo plazo superpuestas en los parámetros & Recuperación más ruidosa y más difícil de rastrear elemento por elemento; no equivale a almacenar el conjunto de datos completo \\
LayerNorm & Estabiliza la norm de los vectores para que la correspondencia distancia/producto punto sea más controlable & Las norms reales y las learned scales siguen modificando la temperature \\
Multi-head & Varios subespacios de representación distintos y distintas retrieval sharpness & No puede equipararse directamente con las microzones de las regiones cerebrales \\
\bottomrule
\end{longtable}

\teachervoice{En la sesión de preguntas, el ponente distingue entre la SDM neuron-centric y la pattern-centric: la attention conserva todas las pattern locations, por lo que la similitud es precisa pero existe una context window; la memoria distribuida de largo plazo solo conserva una noisy superposition, con mayor capacidad pero con dificultad para triangular con precisión el pattern original. Este tradeoff se acerca más a la verdadera enseñanza de ingeniería de la clase que las métricas de monitoreo en producción que inventaba la versión anterior.}

\subsection{Resumen de la sección}

La SDM continua convierte el radio duro en un kernel exponencial y comparte su forma con la attention. Las explicaciones de la learned beta, la LayerNorm y la FFN extienden el marco al Transformer completo, pero cada correspondencia debe tratarse como una hipótesis comprobable. La precisión en contextos cortos de la pattern memory frente a la alta capacidad ruidosa de la distributed memory es el compromiso central que plantea la clase.
\section{Biological Plausibility: ¿puede el cerebelo implementar lectura y escritura al estilo SDM?}

Tras la correspondencia teórica, la primera conferencia lleva la pregunta al nivel de la implementación: si la attention-like retrieval solo requiere activaciones de alta dimensión, un canal de escritura independiente y una pooled readout, ¿ya cuenta el cerebro con circuitos similares? El curso elige el cerebellum, no para afirmar que el cerebelo «es un Transformer», sino para buscar, elemento por elemento, componentes candidatos para address activation, value writing, superposition y majority-like pooling.

\subsection{Circuito abstracto: canal de direcciones y canal de escritura separados}

La diapositiva de Circuitry parte de una abstract neuron. El pattern de entrada determina qué neuronas de dirección hacen firing; otra vía, la value signal, determina qué deben almacenar esas neuronas; en la lectura, la query vuelve a activar las neuronas, cada una emite su value superpuesto y, al final, una pooled unit los agrega. Esta estructura hetero-associative corresponde directamente a la separación key/value de la attention.

\lecturefigure{slide-36-circuitry-sdm.jpg}{Circuito neuronal abstracto que implementa SDM}{Diapositiva V036 recuperada de la grabación oficial; 00:26:38}

\lecturefigure{slide-37-circuitry-mapping-final.jpg}{Estado final de la correspondencia del circuito abstracto con los cerebellar components}{Diapositiva V037 recuperada de la grabación oficial; 00:28:22}

\begin{knowledgebox}{Términos clave: papeles candidatos en el cerebellar circuit}
Las \term{mossy fibers} transportan la entrada; las \term{granule cells} forman una expansión dispersa de alta dimensión; las \term{parallel fibers} envían ampliamente la granule activity a las Purkinje cells; las \term{climbing fibers} aportan una señal de enseñanza/escritura intensa y relativamente independiente; las \term{Purkinje cells} agregan una gran cantidad de entradas e influyen en los deep cerebellar nuclei. El curso alinea estos papeles en el nivel funcional.
\end{knowledgebox}

\subsection{Anatomía del cerebelo y correspondencia, operación por operación, con SDM}

La diapositiva de anatomía muestra los microcircuitos repetidos y relativamente uniformes del cerebelo; la diapositiva final de mapping superpone las vías abstractas de read/write sobre cell types concretos. Las granule cells son enormemente numerosas y tienen conexiones dispersas, por lo que pueden funcionar como hard locations; el mossy input determina la activación; la señal intensa de la climbing fiber puede modificar una Purkinje synapse específica; la Purkinje cell recibe una cantidad enorme de parallel-fiber input y proporciona la pooled output.

\lecturefigure{slide-38-cerebellum-anatomy.jpg}{Tipos de neuronas y estructura de conexiones del cerebellum}{Diapositiva V038 recuperada de la grabación oficial; 00:31:04}

\lecturefigure{slide-39-sdm-cerebellum-mapping.jpg}{Comparación completa entre read/write de SDM y la cerebellar circuitry}{Diapositiva V039 recuperada de la grabación oficial; 00:34:00}

\teachervoice{El ponente afirma que alrededor del 70\% de las neuronas del cerebro se ubican en el cerebellum y advierte que posiblemente no solo participa en el control motor fino; sin embargo, también describe la evidencia sobre cognición de alto nivel como un área de investigación activa. Estas notas conservan ese tono exploratorio y no convierten el hecho de que una región «se ilumine» en la fMRI en un módulo de attention lingüística ya establecido.}

\subsection{Cómo pasar de «implementable» a «comprobado»}

La sesión de preguntas del curso marca el límite científico más importante: aunque el circuito pueda implementar SDM, eso no significa que el cerebro realmente lo use. Una verificación auténtica requiere registrar en tiempo real la sparse neuron activation, el synaptic update y el behavior change durante una associative task; idealmente, también escribir nuevas asociaciones y observar cómo se desvanecen las anteriores. El ponente propone la calcium/voltage imaging de alta velocidad como vía candidata; se trata de un diseño experimental, no de una conclusión ya obtenida.

\begin{warningbox}{Tres brechas de evidencia que no pueden omitirse}
Primero, que la conectividad anatómica permita cierto cómputo no significa que la dinámica lo ejecute; segundo, que la lectura y escritura local se parezcan no significa que todo el cerebro use el training objective de un Transformer; tercero, la analogía entre microzone y multi-head es atractiva, pero el ponente afirma explícitamente que no conviene llevarla demasiado lejos.
\end{warningbox}

En la Q\&A también se distingue entre la fast neural activity y la long-term synaptic plasticity, y se discute la compensación entre capacidad y fidelity que implica el radius: con un radio de activación grande, cada lectura agrega más neuronas y la estimación es más suave, pero aumenta la interferencia entre patterns; con un radio pequeño, la capacidad de discriminación es alta, pero una query con ruido puede activar muy pocas posiciones. En la attention, el parámetro correspondiente no es un radio físico fijo, sino la effective sharpness que producen conjuntamente las key/query norms y beta.

\lecturefigure{slide-40-first-talk-thank-you.jpg}{Resumen de la primera conferencia, referencias y agradecimientos}{Diapositiva V040 recuperada de la grabación oficial; 00:35:54}

\begin{importantbox}{Las conclusiones más sólidas y más débiles de la primera conferencia}
La conclusión más sólida es que, bajo condiciones explícitas, existe una aproximación derivable y verificable numéricamente entre la SDM intersection read y la softmax attention. Una conclusión más débil, aunque sugerente, es que el cerebellar circuit cuenta con componentes capaces de implementar operaciones similares de address/value/readout. Ambas deben presentarse por separado: la primera es evidencia matemática; la segunda, una hipótesis de neurociencia.
\end{importantbox}

\subsection{Resumen de la sección}

El mapping del cerebelo muestra una posible implementación biológica de una attention-like associative memory: la expansión dispersa, la señal de enseñanza independiente, el almacenamiento distribuido y la pooled readout tienen correspondencias funcionales. Sin embargo, la teacher voice más importante de la clase es que «implementable no equivale a verificado»: los experimentos deben seguir directamente los cambios de conectividad durante el aprendizaje y sus resultados conductuales.
\section{Transición a la segunda conferencia: de la memoria asociativa al Cognitive Map}

A partir de 00:42:28 cambia el ponente. El invitado previsto canceló a última hora, y Will Dorrell, con muy poco tiempo de preparación, dio la explicación escribiendo a mano en una tableta; por eso el formato visual pasa de un slide deck estándar a una página larga continua. El contenido, sin embargo, está estrechamente ligado a la primera mitad: Bricken pregunta «cómo una query similar recupera una memory», y Dorrell pregunta «cómo se recuerda la estructura de relaciones de la experiencia y cómo se reutiliza rápidamente con contenido nuevo». La primera pregunta enfatiza la geometría de direcciones; la segunda, la factorization entre estructura y contenido.

\subsection{Problema central: ¿puede reutilizarse la estructura de una experiencia secuencial?}

La primera página manuscrita reúne el hippocampus, el entorhinal system y los Transformers. La página siguiente usa un esquema de sequence/graph para ilustrar un problema general: los objetos y las recompensas que observa el agent pueden cambiar, pero la transition structure, la relative position o la relational rule pueden mantenerse sin cambios. Si el modelo memoriza desde cero el episode completo en cada entorno nuevo, no puede hacer few-shot generalize; si logra extraer la structure, puede transferir las reglas anteriores al contenido nuevo.

\lecturefigure{slide-41-cognitive-map-title.jpg}{Segunda conferencia: Hippocampus, entorhinal system y Transformers}{Estado didáctico manuscrito del video oficial V041; 00:43:54}

\lecturefigure{slide-42-one-problem-sequence-structure.jpg}{Un problema común: extraer una estructura reutilizable de la sequence experience}{Estado didáctico manuscrito del video oficial V042; 00:45:36}

\teachervoice{El moderador del curso explica que Will sustituye al invitado de manera provisional y presenta su formación en computational neuroscience en la Gatsby Unit. Acto seguido, Will describe su propio trabajo así: un modelo que originalmente estudiaba de forma independiente el hippocampal--entorhinal system y que después «resultó parecerse un poco a un Transformer». Esta frase enfatiza la convergent computation, no el diseño de un modelo del cerebro a imagen del Transformer.}

\begin{importantbox}{Definición mínima de structural generalization}
Dados entornos con distintas observation labels pero que comparten el mismo transition/relational graph, el sistema debe reutilizar, tras poca experiencia nueva, la estructura que ya tiene para hacer prediction, planning o inference. Lo que se transfiere es «cómo se conectan las relaciones», no la identidad sensorial de cada nodo.
\end{importantbox}

\subsection{Resumen de la sección}

La segunda conferencia lleva el problema de la memoria de la «recuperación por similitud» a la «reutilización de estructura». Un mapa cognitivo no solo debe recordar lo que ocurrió, sino también separar el what del where/how-related, para que los objetos nuevos puedan vincularse con rapidez a un grafo de relaciones conocido.
\section{Cognitive Map: el espacio es solo un caso particular del razonamiento estructural}

El mapa cognitivo surgió originalmente de la navigation, pero el objetivo del curso es extenderlo a cualquier relational structure. Primero revisamos la evidencia espacial, después la transitive inference no espacial y, por último, abstraemos el problema común como graph learning: el contenido observado cambia, las relaciones entre aristas se mantienen y el modelo necesita reutilizar la estructura y actualizar con rapidez la vinculación del contenido.

\subsection{Del mapa de Tolman a la generalización en árbol}

La diapositiva Theory/navigation describe el cognitive map como un modelo interno de la estructura del entorno, que permite al agent replanificar cuando un camino se bloquea o cambia la meta, en lugar de limitarse a reproducir una cached action sequence. El ejemplo en árbol muestra además que la estructura no tiene que ser un plano bidimensional: las jerarquías, las relaciones de parentesco y los estados de una tarea también pueden formar un graph; lo decisivo es si el modelo puede inferir trayectorias que no ha experimentado directamente.

\lecturefigure{slide-43-cognitive-map-theory-navigation.jpg}{Teoría del cognitive map: un modelo estructural permite una navegación flexible}{Estado de enseñanza manuscrito del video oficial V043; 00:47:32}

\lecturefigure{slide-44-tree-structured-generalization.jpg}{Generalización estructural en relaciones en forma de árbol}{Estado de enseñanza manuscrito del video oficial V044; 00:48:20}

Si el entorno se representa como un grafo $G=(V,E)$, cada nodo $i$ tiene una observación $x_i$ y el tipo de arista o acción $a_t$ determina la transition $i_t\to i_{t+1}$. La generalización estructural exige aprender
$$
g_{t+1}=f_{\theta}(g_t,a_t),
$$
donde $g_t$ es un structural state lo más desacoplado posible de la observation identity concreta. Al cambiar el conjunto de $x_i$, $f_{\theta}$ sigue pudiendo actualizar la posición o las relaciones; solo hace falta volver a vincular $g_i$ con el nuevo contenido.

\begin{knowledgebox}{Concepto previo: un cognitive map no equivale a un mapa geográfico}
Lo esencial de un mapa es que las relaciones se pueden componer: si A está a la izquierda de B y B a la izquierda de C, se deduce la posición relativa de A y C; si cierta acción lleva del nodo 1 al nodo 2, la misma regla de acción puede transferirse a objetos nuevos. El espacio bidimensional es el caso más intuitivo, pero los graph, los family tree, los estados de una tarea y las relaciones entre conceptos pueden usar la misma idea estructural.
\end{knowledgebox}

\subsection{Place cells, grid cells y evidencia de estructura espacial}

La diapositiva Evidence overview reúne tareas espaciales, neural recording y resultados conductuales. El ponente repasa de palabra la place cell: cierta hippocampal neuron presenta un firing intenso cuando el animal se encuentra en una ubicación específica del entorno; la grid cell, en cambio, forma una hexagonal lattice sobre múltiples ubicaciones. No son coordenadas GPS directas, sino que la actividad neuronal muestra una estructura estable en función de la posición.

\lecturefigure{slide-45-task-evidence-selection.jpg}{Panorama de la evidencia conductual y neuronal del mapa cognitivo}{Estado de enseñanza manuscrito del video oficial V045; 00:51:44}

\teachervoice{Will usa el examen «Knowledge» de los London taxi drivers para ilustrar el costo conductual de aprender una estructura espacial, y cita la evidencia de correlación entre el hippocampal size y la experiencia de manejo. Enseguida pasa a tareas no espaciales, porque una «correlación espacial» todavía no demuestra que la región cerebral represente una relational structure más general.}

\begin{warningbox}{Interpretaciones erróneas frecuentes de las tuning curve neuronales}
Que una cell presente firing en una ubicación específica no significa que ella sola «almacene esa ubicación»; la posición suele representarse mediante un population code. Que el grid pattern se parezca a la position encoding de un modelo tampoco implica que el proceso de entrenamiento, el origen de las entradas y la función biológica sean idénticos. Lo que debe compararse es la información decodificable, las regularidades de transformación y el papel causal en la tarea.
\end{warningbox}

\subsection{Evidencia no espacial: Transitive inference}

La diapositiva Transitive-inference presenta un tipo de tarea relacional no espacial: el animal aprende pair locales como A mayor que B, B mayor que C, etc., y después se le evalúa con A y C, que no se entrenaron directamente. Si solo memoriza los observed pairs, no puede inferir de forma estable las combinaciones nuevas; si forma un ordered relational map, puede deducirlas siguiendo la estructura. La diapositiva manuscrita coloca los experimentos conductuales, la evidencia de imagen y la pair relation en una misma cadena argumentativa.

\lecturefigure{slide-46-transitive-inference-evidence.jpg}{Transitive inference: inferir relaciones no vistas a partir de pair aprendidos}{Estado de enseñanza manuscrito del video oficial V046; 00:54:54}

\lecturefigure{slide-47-executive-layer-map.jpg}{Observation, structure y lectura conductual en tareas relacionales}{Estado de enseñanza manuscrito del video oficial V047; 00:55:18}

\begin{knowledgebox}{Lectura de la figura: cómo distinguir las tres estrategias}
\term{pair memorization} solo puede responder sobre los pair entrenados; \term{scalar ordering} asigna a cada objeto un rank unidimensional, adecuado para órdenes lineales; \term{relational graph} representa varios tipos de edge y puede manejar árboles, ciclos y espacios bidimensionales. El diseño experimental debe distinguir estas estrategias mediante combinaciones no vistas y cambios de estructura, en lugar de medir solo la pair accuracy en entrenamiento.
\end{knowledgebox}

\subsection{Del graph world model a un algoritmo neuronal general}

La diapositiva Graph-world abstrae el entorno como nodos, aristas y observation; la siguiente enumera el algoritmo objetivo: la estructura debe poder reutilizarse entre entornos, el contenido debe escribirse con rapidez, la planificación debe ocurrir sobre el latent graph y las representaciones neuronales deben contrastarse con datos hippocampal/entorhinal. El reto abarca a la vez representation learning, memory binding y model-based inference.

\lecturefigure{slide-48-graph-world-model.jpg}{Abstraer el entorno como un graph reutilizable y observations variables}{Estado de enseñanza manuscrito del video oficial V048; 00:58:58}

\lecturefigure{slide-49-neural-algorithm-general.jpg}{Algoritmo neuronal ideal: transferencia estructural, vinculación rápida e inferencia}{Estado de enseñanza manuscrito del video oficial V049; 00:59:28}

\begin{importantbox}{Descomposición de la tarea en la segunda clase}
El aprendizaje lento se encarga de dominar la transition structure y las reglas componibles; el aprendizaje rápido se encarga de vincular el nuevo sensory content con los structural states; la inferencia se encarga de propagar y planificar a lo largo de la estructura. La arquitectura de TEM asigna precisamente estas tres funciones a representaciones y mecanismos de actualización distintos.
\end{importantbox}

\subsection{Resumen de la sección}

Las place/grid cells y la taxi-driver evidence respaldan el mapa cognitivo espacial, mientras que tareas como la transitive inference respaldan un relational map más amplio. El problema computacional común es separar la observation identity de la reusable structure y después vincularlas mediante memoria rápida, de modo que también puedan inferirse combinaciones no vistas.
\section{Tolman--Eichenbaum Machine: What, Where y Associative Memory}

TEM es el núcleo del modelo de la segunda parte de la clase. Hace una analogía entre la lateral entorhinal cortex (LEC) y el content/what stream, entre la medial entorhinal cortex (MEC) y el structural/location stream, y entre el hippocampus y la conjunctive associative memory. El estado estructural se actualiza de forma recurrente (recurrently) según la action, el contenido se vincula rápidamente en un entorno nuevo, y la memory sostiene después el recall y la inference.

\subsection{Panorama de la arquitectura: cada uno de los tres estados cumple su función}

Después de la portada, el diagrama de la arquitectura muestra tres capas de información: sensory observation, structural state reutilizable y conjunctive memory. El graph de la parte inferior indica que una misma estructura puede alojar distintas observation labels. El modelo no mete la observation directamente en una RNN, sino que hace que la structure dynamics y el content encoding se aprendan por separado y luego se combinen en un hippocampal-like state.

\lecturefigure{slide-50-tem-title.jpg}{Tolman--Eichenbaum Machine: unificación de la memoria espacial y la relacional}{Estado didáctico manuscrito del video oficial V050; 00:59:38}

\lecturefigure{slide-51-tem-architecture.jpg}{Arquitectura de TEM: sensory content, structural state y memory}{Estado didáctico manuscrito del video oficial V051; 01:01:08}

\begin{longtable}{@{}L{0.17\textwidth}L{0.31\textwidth}L{0.40\textwidth}@{}}
\toprule
Componente & Función computacional & Analogía neurocientífica \\
\midrule
$x_t$ / sensory code & Representación del contenido de la observation actual & LEC-like what/content stream \\
$g_t$ / structural code & Representación de posición/relación actualizada a partir de la acción y el estado anterior & MEC-like where/structure stream \\
$p_t$ / conjunctive code & Vincula el contenido actual con la estructura en un episode & hippocampal conjunctive representation \\
$M_t$ / associative memory & Escribe rápidamente los conjunctive states y los recupera a partir de un cue & hippocampal recurrent/associative memory \\
\bottomrule
\end{longtable}

\subsection{LEC y MEC: separación entre contenido y estructura}

La diapositiva de LEC extrae el sensory code a partir de la observation y responde «qué se vio»; la diapositiva de MEC actualiza el structural code según la action/transition y responde «en qué lugar de la estructura se está». Esta separación permite que el modelo conserve la transition dynamics en un entorno nuevo y, al mismo tiempo, sustituya la object identity. Si ambas estuvieran completamente entrelazadas desde el inicio, cada cambio de observation podría destruir el mapa existente.

\lecturefigure{slide-52-lec-what-stream.jpg}{LEC-like stream: codificación del sensory content}{Estado didáctico manuscrito del video oficial V052; 01:01:20}

\lecturefigure{slide-53-mec-where-stream.jpg}{MEC-like stream: actualización del reusable structural state}{Estado didáctico manuscrito del video oficial V053; 01:02:24}

La actualización estructural puede abstraerse como
$$
g_t=f_{\theta}(g_{t-1},a_t),
\qquad
x_t=e_{\phi}(o_t),
$$
donde $o_t$ es la observation, $a_t$ es la action/edge, $g_t$ es el path-integrated structural state y $x_t$ es el content encoding. $\theta$ se aprende lentamente a través de los entornos, mientras que el object binding del entorno actual puede establecerse con rapidez.

\begin{warningbox}{«Where» no es necesariamente una coordenada bidimensional}
En una tarea de tree o de transitive-order, $g_t$ representa una relational position y no tiene que corresponder a un $(x,y)$ euclidiano. Interpretar el MEC-like state como un código estructural general es justamente lo que permite a TEM generalizar de la navigation a la relational memory; sin embargo, esa generalización todavía requiere el respaldo de tareas y de datos neuronales.
\end{warningbox}

\subsection{Hippocampal binding y Hebbian-style memory}

La diapositiva de memory-update combina content y structure en un conjunctive code, que después se escribe en la associative memory. La diapositiva Hebbian ofrece la intuición de una actualización local en la que «la activación conjunta refuerza la conexión». Una forma abstracta de escribirlo es
$$
p_t=b(x_t,g_t),
\qquad
M_t=\lambda M_{t-1}+p_t p_t^{\top},
$$
donde $b$ es la binding operation, $p_t$ es la conjunctive representation, $M_t$ es la memory matrix y $\lambda$ controla cuánto se conserva la memoria antigua. El TEM real usa una estructura de inference/generative más detallada, pero esta forma simplificada muestra por qué la escritura rápida no exige volver a entrenar todos los parámetros.

\lecturefigure{slide-54-recurrent-memory-update.jpg}{Recurrent state y associative-memory update en TEM}{Estado didáctico manuscrito del video oficial V054; 01:02:58}

\lecturefigure{slide-55-hebbian-memory-equations.jpg}{Hebbian-style binding y memory equations}{Estado didáctico manuscrito del video oficial V055; 01:03:40}

\begin{knowledgebox}{Qué son los fast weights}
Los parámetros del modelo $\theta,\phi$ se entrenan lentamente a lo largo de muchos entornos, mientras que la memory matrix $M_t$ se actualiza rápidamente dentro del episode actual; por eso suele llamársele fast weights. Esto permite guardar contenido nuevo tras una sola exposure o unas pocas, mientras la dynamics estructural sigue proviniendo de los parámetros de largo plazo.
\end{knowledgebox}

\subsection{Graph inference y planning}

La diapositiva de graph-inference muestra cómo el agent usa el structural state y la recalled memory para inferir nodos que no observó directamente o para planear rutas. Lo esencial no es almacenar el graph completo como una adjacency table fija, sino hacer que el recurrent transition model genere una structure candidata y que después la memory restituya la observation identity del entorno actual.

\lecturefigure{slide-56-graph-inference-plan.jpg}{Graph inference a partir del structural state y la memory}{Estado didáctico manuscrito del video oficial V056; 01:06:02}

\begin{lstlisting}[caption={Flujo conceptual del estado estructural y la vinculación rápida}]
def tem_step(structure, action, observation, memory):
    next_structure = transition_model(structure, action)
    content = content_encoder(observation)
    conjunctive = bind(next_structure, content)
    memory.write(conjunctive)
    recalled = memory.read(next_structure)
    return next_structure, recalled
\end{lstlisting}

\teachervoice{Will describe el objetivo como «la estructura se aprende lentamente; el contenido se vincula rápidamente». Esto hace eco de la separación address/value de SDM en la primera mitad: la structure es como un address system transferible, el sensory content es como el value que se escribe, y la hippocampal memory combina ambos de forma temporal.}

\subsection{Grid-like code y adaptación rápida}

Las dos diapositivas de grid-code comparan la representation del modelo con el periodic spatial pattern que suele aparecer en los datos neuronales; la diapositiva de learning curve muestra que el agent mejora rápidamente después de pocas visitas. La conclusión correcta es que un recurrent state estructurado, junto con el entrenamiento en la tarea, puede producir una grid-like basis y sostener un map learning rápido. Eso no significa que cada hidden unit corresponda a una grid cell real, ni demuestra que la retícula hexagonal sea la única codificación viable.

\lecturefigure{slide-57-grid-code-evidence.jpg}{Grid-like representations que surgen en TEM}{Estado didáctico manuscrito del video oficial V057; 01:08:26}

\lecturefigure{slide-58-grid-code-summary.jpg}{Comparación estructural entre el modelo y los neural grid codes}{Estado didáctico manuscrito del video oficial V058; 01:09:36}

\lecturefigure{slide-59-quick-learning-result.jpg}{Curva de aprendizaje que mejora rápidamente tras visitar pocos nodos}{Estado didáctico manuscrito del video oficial V059; 01:10:48}

\begin{knowledgebox}{Lectura de la figura: grid-like no se juzga solo por «si parece hexagonal»}
Hay que examinar la periodicidad, la orientación, la escala y la fase del firing field, así como su estabilidad ante cambios del entorno, y comparar la population structure. Las manchas repetidas en un solo mapa de calor pueden deberse al suavizado, a los bordes o a los datos de entrenamiento; solo cuando varias estadísticas geométricas coinciden con la causalidad de la tarea se tiene una evidencia más sólida.
\end{knowledgebox}

\subsection{Resumen de la sección}

TEM logra un aprendizaje rápido de mapas mediante un content stream, un structural recurrent state y una fast associative memory. La representación estructural se reutiliza entre entornos, las observation nuevas se vinculan rápidamente y la memory sostiene la graph inference. El grid-like code y el few-visit learning son evidencia funcional, pero no deben presentarse como una identidad anatómica.
\section{Por qué TEM «parece un Transformer»}

Al final de la segunda clase, la hippocampal associative memory se reformula como una modern Hopfield network. La modern Hopfield update y la softmax attention comparten la similitud exponencial y la ponderación por value; por eso, la memory read de TEM puede expresarse directamente con query, keys y values. Esta reformulation no es solo una analogía: también mejora la expansión de dimensión del outer-product binding y permite añadir context adicional como una entrada independiente.

\subsection{De la TEM equation a la Attention update}

La primera página manuscrita pone lado a lado la recurrent/associative update de TEM y la fórmula del Transformer. Sea $q_t$ la cue actual, $K=[k_1,\ldots,k_t]$ las memory keys históricas y $V=[v_1,\ldots,v_t]$ los values; entonces, la lectura de la memoria moderna es
$$
m_t=V\operatorname{softmax}(\beta K^{\top}q_t).
$$
Aquí, $q_t$ es la hippocampal/structural cue actual, $k_\tau$ es una memory address pasada y $v_\tau$ es el bound content que se debe recuperar. Si solo se permite $\tau\le t$, también cumple la restricción temporal de la causal attention.

\lecturefigure{slide-60-tem-transformer-equation.jpg}{Reformulación de la TEM memory update como una Transformer-like equation}{Estado manuscrito de la clase en el video oficial V060; 01:11:52}

\lecturefigure{slide-61-attention-query-retrieval.jpg}{La associative recall desde la perspectiva de Query, keys y values}{Estado manuscrito de la clase en el video oficial V061; 01:12:34}

\begin{importantbox}{El cálculo común de Modern Hopfield y Attention}
El neural state actual es la query; las memories almacenadas aportan las keys y los values; el dot product mide la similitud, y la softmax selecciona y mezcla las memories. Las diferencias provienen más de cómo se generan las representaciones, del ciclo de vida de la memory, del objetivo de entrenamiento y de la implementación biológica que de la forma superficial de la read equation.
\end{importantbox}

\subsection{Pasos de actualización de la Relational memory}

La página manuscrita larga descompone el proceso en la generación de key/query, la value encoding, la causal mask, la memory retrieval y la recurrent state update. Aquí, la positional encoding no es una sinusoid fija, sino el structural state que actualiza el transition model; este incorpora el inductive bias del graph del entorno, por lo que una sola attention layer puede bastar para el razonamiento relacional que un Transformer común necesita aprender con varias capas.

\lecturefigure{slide-62-relational-memory-steps.jpg}{key/query/value y recurrent update de la Relational memory}{Estado manuscrito de la clase en el video oficial V062; 01:14:56}

\begin{knowledgebox}{Términos clave: la positional encoding es distinta en los dos sistemas}
El Transformer estándar suele codificar el sequence index como posición; el $g_t$ de TEM se parece más a una latent graph position obtenida por la acumulación de actions. En ambos casos, un mismo content produce representaciones distintas en posiciones estructurales distintas, pero en TEM la actualización de la posición está restringida explícitamente por la transition dynamics, por lo que su sesgo inductivo es más fuerte.
\end{knowledgebox}

\subsection{Hopfield memory clásica y moderna}

La página Neural-memory-softmax repasa lo siguiente: la classic Hopfield recupera el attractor más cercano mediante recurrent energy descent, y su capacidad y su interferencia son limitadas; la modern Hopfield usa similitud exponencial y memory patterns explícitos, su fórmula de actualización se aproxima a la attention y aumenta notablemente la capacidad. El curso aprovecha precisamente esta última, sin afirmar que toda la biological hippocampal recurrence use softmax.

\lecturefigure{slide-63-neural-memory-softmax.jpg}{La softmax retrieval de la Modern Hopfield memory}{Estado manuscrito de la clase en el video oficial V063; 01:15:22}

Puede expresarse con una intuición energética unificada:
$$
E(q)=-\frac{1}{\beta}\log\sum_{\mu}\exp(\beta k_\mu^{\top}q)
+\frac{1}{2}\|q\|^2.
$$
El primer término favorece que $q$ se acerque a ciertos stored patterns, y el segundo restringe la magnitud del estado. Actualizar sobre esta energía produce una softmax-weighted memory combination. La expresión explica por qué la memory capacity y la retrieval sharpness pueden modificarse mediante el kernel exponencial.

\subsection{Correspondencia arquitectónica, beneficios de escalamiento y limitaciones}

La página Comparison alinea LEC, MEC e hippocampus con el content encoder, el recurrent position state y la attention memory; la página analysis continúa con la comparación entre graph transitions y attention paths. La versión Modern Hopfield evita hacer un enorme outer product entre content y position para luego aplicar flatten; además, el context adicional puede añadirse como un nuevo flujo de entrada, y la complejidad pasa de una expansión multiplicativa de dimensiones a una combinación casi lineal de varios vectores.

\lecturefigure{slide-64-tem-transformer-comparison.jpg}{Correspondencia funcional entre TEM y los Transformer-like components}{Estado manuscrito de la clase en el video oficial V064; 01:17:46}

\lecturefigure{slide-65-tem-analysis-diagrams.jpg}{Análisis estructural de la TEM/Transformer correspondence}{Estado manuscrito de la clase en el video oficial V065; 01:18:20}

\teachervoice{La conclusión de Will es muy mesurada: el fuerte inductive bias facilita la tarea, y por eso basta una capa de attention; esto no demuestra que «un Transformer de una capa lo resuelve todo». Él sitúa el valor en la interpretabilidad, las predicciones neuronales, la aceleración del código y el binding escalable de context adicional.}

\begin{warningbox}{Las dos caras de un sesgo inductivo fuerte}
Si la tarea realmente se compone de una graph transition estable y de un content binding rápido, TEM es muy eficiente; si la propia estructura cambia con frecuencia, si las acciones no son observables o si el contenido y la estructura están fuertemente acoplados, la factorization fija puede convertirse en un sesgo erróneo. Al compararlo con un Transformer de propósito general, hay que igualar el alcance de la tarea, no solo comparar el número de capas.
\end{warningbox}

\subsection{Conclusiones y Q\&A sobre Grid-cell}

La página de conclusiones resume los beneficios en ambos sentidos: la modern memory de la IA le da a TEM mejor scaling e implementación; el recurrent position code, la content/structure separation y el grid-like state de la neuroscience aportan hipótesis para el diseño y la explicación de los Transformer. Al final, la Q\&A explica cómo se mide una sola grid cell mediante electrode spike sorting, cómo distintas cells pueden compartir escala/orientación con fases desplazadas y cómo forman varios lattice-scale modules.

\lecturefigure{slide-66-conclusions.jpg}{Conclusiones de la segunda clase: la relación bidireccional entre Neuroscience y AI}{Estado manuscrito de la clase en el video oficial V066; 01:19:42}

\teachervoice{Un estudiante pregunta si la grid template es producto de la evolution o del learning. Will responde que su estructura regular aparece muy temprano en animales jóvenes, lo que muestra un fuerte developmental/evolutionary bias, pero aún no está claro si el sistema puede ser co-opt por otras relational tasks ni qué tan flexible es. Esta incertidumbre tiene más valor didáctico que la frase «las grid cells son una codificación posicional».}

\subsection{Resumen de la sección}

La Modern Hopfield retrieval expresa la TEM memory como attention: el state actual es la query, las conjunctive memories históricas aportan keys/values y la softmax realiza la recuperación. El recurrent structural code se asemeja a una positional encoding adaptada a la tarea. La similitud ofrece un puente interpretable y beneficios de implementación, pero el fuerte sesgo estructural de TEM, el alcance limitado de sus tareas y sus hipótesis neuronales restringen la extrapolación.
\section{Síntesis y ampliación}

Las dos conferencias explican la Transformer-like computation desde direcciones distintas. La ruta de SDM parte de la geometría de alta dimensión: el intersection count de las distributed addresses aproxima una similitud exponencial, de donde se obtiene la softmax retrieval; el cerebellar circuit ofrece una implementación candidata de read/write. La ruta de TEM parte de las necesidades cognitivas: separa la reusable structure del changing content, los vincula mediante una fast associative memory y luego reformula la memory read como modern Hopfield attention.

\subsection{Comparación unificada de las dos rutas}

La siguiente tabla no fusiona el cerebellum y el hippocampus en un solo modelo cerebral; compara qué parte de la attention explica cada uno. La primera conferencia se centra más en el weighting kernel y en la geometría del almacenamiento; la segunda, en el positional/structural state, el rapid binding y la compositional inference. Ambas enfatizan la key/content separation, la distributed memory y la uncertainty about biological implementation.

\begin{longtable}{@{}L{0.18\textwidth}L{0.24\textwidth}L{0.25\textwidth}L{0.21\textwidth}@{}}
\toprule
Ruta & Pregunta principal & Correspondencia en el Transformer & Evidencia más sólida \\
\midrule
SDM geometry & Cómo una noisy query recupera recuerdos según su similitud & softmax attention weights & Aproximación de intersecciones en alta dimensión y experimentos numéricos \\
Cerebellar mapping & Cómo un circuito neuronal separa activación, escritura y lectura & keys/values y pooled retrieval & Factibilidad anatómica y experimentos candidatos \\
Cognitive map / TEM & Cómo reutilizar la estructura relacional y vincular rápidamente contenido nuevo & recurrent position state + memory & Comparación entre conducta, modelos y representaciones neuronales \\
Modern Hopfield TEM & Cómo ampliar la capacidad de la associative read e integrarla con el context & query/key/value attention & Correspondence de fórmulas y beneficios de implementación \\
\bottomrule
\end{longtable}

\begin{importantbox}{El verdadero método «neuroscience-inspired» del curso}
No se trata de copiar el nombre de alguna región cerebral, sino de identificar primero el problema computacional, proponer después un modelo con un inductive bias explícito y, por último, establecer vínculos falsables entre las matemáticas, la conducta en la tarea, las representaciones internas, los datos neuronales y los experimentos de intervención. Una brain analogy vale más que una metáfora solo cuando genera predicciones nuevas.
\end{importantbox}

\subsection{Seis lecciones para la investigación sobre el Transformer}

\begin{enumerate}
  \item \textbf{Explicar las condiciones geométricas del softmax.} Revisar la norm, la temperature, la dimension y la approximation region, en lugar de tratar el softmax como un valor predeterminado misterioso.
  \item \textbf{Distinguir la pattern memory de corto plazo de la distributed memory de largo plazo.} La context precision y la parameter capacity enfrentan interferencias y grados de trazabilidad distintos.
  \item \textbf{Dejar explícita la separación entre key y value.} La señal de coincidencia y el contenido devuelto son distintos; este es el mecanismo común a la attention, la SDM y el hippocampal binding.
  \item \textbf{Hacer que la representación de la posición refleje la estructura de la tarea.} El sequence index es solo la estructura más simple; la graph transition, la action y el relational state pueden aportar positional dynamics más potentes.
  \item \textbf{Reservar un fast state para el vínculo rápido.} El contenido de un entorno nuevo no siempre debe depender de actualizaciones lentas de parámetros; la memory/fast weights puede encargarse del episode-specific knowledge.
  \item \textbf{Convertir la biological plausibility en experimentos.} Especificar qué cells registrar, cuándo se escribe, cómo intervenir y qué resultado refutaría el mapping.
\end{enumerate}

\begin{warningbox}{Las tres afirmaciones de esta clase que más fácilmente se exageran}
«Attention approximates SDM» no significa que la attention y la SDM sean idénticas con cualquier parámetro; «el cerebelo puede implementar la SDM» no significa que esté demostrado que el cerebelo ejecuta un Transformer; «TEM equivale aproximadamente a un Transformer» no significa que el hippocampus sea un modelo de lenguaje. Cada afirmación debe acompañarse del nivel de comparación, los supuestos y el espacio de contraejemplos.
\end{warningbox}

\subsection{Lecturas adicionales}

\begin{itemize}
  \item Bricken and Pehlevan, \href{https://proceedings.neurips.cc/paper/2021/hash/8171ac2c5544a5cb54ac0f38bf477af4-Abstract.html}{\emph{Attention Approximates Sparse Distributed Memory} (NeurIPS)}
  \item Whittington et al., \href{https://www.sciencedirect.com/science/article/pii/S009286742031388X}{\emph{The Tolman--Eichenbaum Machine} (Cell)}
  \item Whittington et al., \href{https://arxiv.org/abs/2112.04035}{\emph{Relating Transformers to Models and Neural Representations of the Hippocampal Formation} (arXiv)}
  \item Bricken et al., \href{https://openreview.net/forum?id=JknGeelZJpHP}{\emph{Sparse Distributed Memory is a Continual Learner} (OpenReview)}
  \item Behrens et al., \href{https://www.nature.com/articles/s41593-022-01153-y}{\emph{How to Build a Cognitive Map} (Nature Neuroscience)}
\end{itemize}

Al leer, conviene aplicar siempre la misma lista de verificación: cuál es el objeto computacional, cómo se factoriza la representación, cómo se escribe y se lee la memory, qué equivalencias son exact/approximate, qué resultados provienen del modelo y cuáles de los neural data, y qué experimento podría refutar el mapping. Solo así la neuroscience se convierte en diseño de modelos y en predicciones científicas, en lugar de limitarse a ponerle etiquetas de regiones cerebrales a arquitecturas existentes.

\end{document}
